\documentclass[10pt]{article}

\usepackage[margin=0.75in]{geometry}
\usepackage{amsmath,amsthm,amssymb,bm}
\usepackage{xcolor}
\usepackage{cancel}
\usepackage{graphicx}
\usepackage{changepage}
\usepackage{circuitikz}
\usepackage{pgfplots}
\usepackage{minted}
\usepackage{physics}
\usepackage{tikz}
\usepackage{tikz-3dplot}
\usepackage{hyperref}
\usepackage{siunitx}
\usepackage{fontspec}
\usepackage{relsize}
\usepackage{subfig}
\usepackage{todonotes}
\usepackage{contour}
\usepackage{multicol, multirow, booktabs}
\usepackage[breakable]{tcolorbox}
\usepackage[inline]{enumitem}

\theoremstyle{definition}
\newtheorem{problem}{Problem}
\newtheorem{soln}{Solution}

\pgfplotsset{compat=newest}
\usetikzlibrary{lindenmayersystems}
\usetikzlibrary{arrows}
\usetikzlibrary{calc}
\usetikzlibrary{positioning, fit}
\usetikzlibrary{3d, perspective}
\usetikzlibrary{math}
\usetikzlibrary{fadings}
\usetikzlibrary{angles,quotes} 
\usetikzlibrary{decorations.markings}

\definecolor{incolor}{HTML}{303F9F}
\definecolor{outcolor}{HTML}{D84315}
\definecolor{cellborder}{HTML}{CFCFCF}
\definecolor{cellbackground}{HTML}{F7F7F7}
\newcommand{\ui}{\hat{i}}
\newcommand{\uj}{\hat{j}}
\newcommand{\uk}{\hat{k}}
\newcommand{\ux}{\hat{\mathbf{x}}}
\newcommand{\uy}{\hat{\mathbf{y}}}
\newcommand{\uz}{\hat{\mathbf{z}}}
\newcommand{\uv}[1]{\hat{\bv{{#1}}}}
\newcommand{\pr}[1]{{#1^\prime}}
\newcommand{\pd}[2][]{{\frac{\partial^{#1}}{\partial {{#2}^{#1}}}}}
\newcommand{\dif}[2][]{{\frac{d^{#1}}{d{{#2}^{#1}}}}}
\newcommand{\grd}{\nabla}

\pgfdeclarelayer{background}  
\pgfsetlayers{background,main}
\AtBeginDocument{\RenewCommandCopy\qty\SI}

\makeatletter
\newcommand{\boxspacing}{\kern\kvtcb@left@rule\kern\kvtcb@boxsep}
\makeatother
\newcommand{\prompt}[4]{
    \ttfamily\llap{{\color{#2}[#3]:\hspace{3pt}#4}}\vspace{-\baselineskip}
}

\newcommand{\thevenin}[2]{
  \begin{center}
    \begin{circuitikz} \draw
      (0,0) -- (2,0) to[battery1, l_=$V_{Th}\eq#1$] (2,2) 
      to[resistor, l_=$R_{Th}\eq#2$] (0,2)
      ;
      \draw [o-] (-.07,2.079);
      \draw [o-] (-.07,0.079);
    \end{circuitikz}
  \end{center}
}

\newcommand{\norton}[2]{
  \begin{center}
    \begin{circuitikz} \draw
      (0,0) -- (3,0) to[american current source, l_=$I_{N}\eq#1$] (3,2) -- (0,2) (2,0)
      to[resistor, l=$R_{N}\eq#2$] (2,2)
      ;
      \draw [o-] (-.07,2.079);
      \draw [o-] (-.07,0.079);
    \end{circuitikz}
  \end{center}
}

\DeclareMathOperator{\Div}{div}
\DeclareMathOperator{\Curl}{curl}
\DeclareMathOperator{\sgn}{sgn}

\newcommand{\highlight}[1]{\colorbox{yellow}{$\displaystyle #1$}}
\newcommand{\ti}[1]{\widetilde{#1}}
\renewcommand{\div}[1]{\bv{\nabla}\cdot #1}
\renewcommand{\curl}[1]{\bv{\nabla}\times #1}

\newfontface{\Kaufmann}{Kaufmann}
\DeclareTextFontCommand{\kf}{\Kaufmann}
\newcommand{\scriptr}{\fontsize{12pt}{12pt}\kf{r}\fontsize{10pt}{10pt}}

\newfontface{\KaufmannB}{Kaufmann Bold}
\DeclareTextFontCommand{\kfb}{\KaufmannB}
\newcommand{\bscriptr}{\fontsize{12pt}{12pt}\kfb{r}\fontsize{10pt}{10pt}}
\newcommand{\justif}[2]{&{#1}&\text{#2}}
\newcommand{\bv}[1]{\bm{\mathrm{#1}}}

\title{Math 3790H: Assignment II}
\author{Jeremy Favro (0805980) \\ Trent University, Peterborough, ON, Canada}
\date{\today}

\begin{document}
\maketitle

% PROBLEM 1
\begin{problem}
Let $s_1 = 1$ and $s_{n+1} = \frac{1}{3} (s_n + 1)$ for $n \geq 1$.
\begin{enumerate}[label=(\alph*)]
  \item Find $s_2$, $s_3$ and $s_4$.
  \item Use induction to show $s_n > 1/2$ for all $n$. [4]
  \item Show that $(s_n)$ is a decreasing sequence. [3]
  \item Show lim $s_n$ exists and find lim $s_n$. [5]
\end{enumerate}
\end{problem}
\begin{soln}

\end{soln}

% PROBLEM 2
\begin{problem}
Let $r\in \mathbb{R}$ with $0 < r < 1$. Suppose $(s_n)$ is a sequence having the property
that
$$|s_{n+2} - s_{n+1}| \leq r|s_{n+1} - s_n|$$
for all $n \in\mathbb{N}$.
\begin{enumerate}[label=(\alph*)]
  \item Show that $(s_n)$ is a Cauchy sequence.
  \item Give an example showing that this can fail for $r = 1$.
\end{enumerate}
\end{problem}
\begin{soln}~
  \begin{enumerate}[label=(\alph*)]
    \item Consider that
          $$|s_{3} - s_{2}| \leq r|s_2 - s_1|$$
          then, repeatedly chaining the condition given
          $$|s_{k+2} - s_{k+1}| \leq r|s_{k+1} - s_k|\leq r^{2}|s_{k} - s_{k-1}|\leq\dots \leq r^{k-1}\abs{s_2 - s_1}.$$
          Consider then $m,n\in\mathbb{N}$ with (WLOG) $m>n$. Then
          $$\abs{s_m-s_n}=\abs{s_m-s_{m-1}+s_{m-1}-s_{m-2}+\dots+ s_{n+1}-s_n}$$
          and by triangle inequality
          $$\abs{s_m-s_n}\leq \abs{s_m-s_{m-1}}+\abs{s_{m-1}-s_{m-2}}+\dots+\abs{s_{n+1}-s_n}.$$
          Then we have that $\abs{s_m-s_{m-1}}\leq r^{m-2}\abs{s_2 - s_1}$, $\abs{s_{m-1}-s_{m-2}}\leq r^{m-3}\abs{s_2 - s_1}$,
          etc. up to $\abs{s_{n+1}-s_n}\leq r^{n-1}\abs{s_2 - s_1}$. Certainly then $\abs{s_m-s_n}$ is bounded above by the sum of all of
          these bounds,
          $$\abs{s_m-s_n}\leq \left[r^{m-2}+r^{m-3}+\dots+r^{n-1}\right]\abs{s_2 - s_1}.$$
          If we factor out an $r^{n-1}$, and reverse the order so it's more clear,
          $$\abs{s_m-s_n}\leq r^{n-1}\left[1+r+r^2+\dots+r^{m-n-2}+r^{m-n-1}\right]\abs{s_2 - s_1}$$
          which is a geometric series from $0$ to $m-n-1$. Such a geometric series is, for the same $r$, bounded above by its infinite form,
          $$\sum_{i=0}^{\infty}r^i=\frac{1}{1-r}$$
          which tells us that 
          $$\abs{s_m-s_n}\leq \frac{r^{n-1}}{1-r}\abs{s_2 - s_1}.$$
          Now the definition of a Cauchy sequence is that for each $\epsilon>0$ (I still don't quite see the difference between for each and for all in this case, but whatever)
          there exists an $N\in\mathbb{N}$ such that $m,n> N$ implies that 
          $$\abs{s_m-s_n}< \epsilon.$$
          Certainly if 
          $$\abs{s_m-s_n}\leq \frac{r^{n-1}}{1-r}\abs{s_2 - s_1},$$
          then $N$, which is smaller than $n$, is also suitable for an upper bound as $r^{N-1}>r^{n-1}$ so 
          $$\abs{s_m-s_n}<\frac{r^{N-1}}{1-r}\abs{s_2 - s_1}$$
          by making $N$ arbitrarily large, we can then force $\abs{s_m-s_n}$ arbitrarily small, meaning $s_n$ is Cauchy?
    \item Give an example showing that this can fail for $r = 1$.
  \end{enumerate}
\end{soln}
\end{document}